% !TEX program = xelatex
% 编译方式：在本目录执行两遍 xelatex（详见 docs/_manual_common/README.md）。
% 源码依据：src/carbon_analysis/ 与 include/hacdcpf/carbon_analysis/ 源码及
%           docs/theory/carbon_analysis_contract.md，公式与参数以代码为准。
\documentclass[UTF8,fontset=fandol,a4paper,11pt]{ctexart}
\usepackage{../../_manual_common/hysim_manual}

\title{\textbf{交直流混合高弹性能源电力系统仿真平台}\\[0.5em]
       {\Large 碳流追踪与碳排放核算技术手册}\\[0.3em]
       {\large 数学模型、接口参数、验证校核与工业级评价}}
\author{HySim-XJTU-HRPES（西安交通大学 HRPES 团队）}
\date{\today}

\begin{document}
\maketitle
\begin{abstract}
\noindent 本手册依据 HySim-XJTU-HRPES 平台碳流分析模块
（\srcpath{src/carbon\_analysis/}、\srcpath{include/hacdcpf/carbon\_analysis/}）整理，
覆盖两类基于已收敛潮流结果的后处理方法：比例功率追踪法（BFS 碳流追踪）与
碳排放流矩阵法，以及年度碳核算与用户/节点绿证（GEC）核算。全部结论以源码与
既有契约 \srcpath{docs/theory/carbon\_analysis\_contract.md} 为准；碳流分析为纯后处理，
不改变调度、不修复失败潮流。文档版式与《潮流计算技术手册》《最优潮流技术手册》一致。
\end{abstract}

\tableofcontents
\newpage

\section{阅读指南与约定}
\subsection{数据来源}
\begin{itemize}
  \item \srcpath{include/hacdcpf/carbon\_analysis/carbon\_analysis.hpp}、
        \srcpath{annual\_carbon\_analysis.hpp} —— 快照/年度碳分析接口与结果结构；
  \item \srcpath{src/carbon\_analysis/carbon\_analysis.cpp}、
        \srcpath{annual\_carbon\_analysis.cpp} —— 追踪法与矩阵法实现；
  \item \srcpath{docs/theory/carbon\_analysis\_contract.md} —— 输入/单位/源 ID/交直流身份与
        有效性门控契约；
  \item \srcpath{tests/test\_carbonflow\_tracing.cpp} 等 —— 验证依据。
\end{itemize}
命名空间为 \texttt{hacdcpf::analysis}。

\subsection{单位制与索引空间}
排放量单位 tCO\textsubscript{2}（年度量为 tCO\textsubscript{2}/yr），碳强度单位
tCO\textsubscript{2}/MWh，功率 MW。AC 与 DC 集合分开核算：结果向量分别对应
\fld{load\_carbon}/\fld{dc\_load\_carbon}、\fld{branch\_carbon}/\fld{dc\_branch\_carbon}、
\fld{bus\_carbon}/\fld{dc\_bus\_carbon}，索引空间与对应潮流结果一致。

\subsection{诚实口径}
碳流分析不含单一“成功”布尔量；调用方须检查 \fld{power\_balance\_verified}、
\fld{matrix\_solved}/\fld{matrix\_rank}/\fld{matrix\_relative\_residual}、
\fld{tracing\_verified}。潮流未收敛时输出默认“未验证”结果；奇异/病态矩阵保持可见
（不会被伪装为零碳）。\fld{loss\_allocation\_alpha} 只重新分配损耗排放的归属责任，
不改变物理损耗。

\section{功能定位与架构}
\subsection{公共入口族}
\begin{center}\footnotesize
\begin{tabular}{@{}L{6.6cm}L{3.0cm}L{5.4cm}@{}}
\toprule
入口 & 定义位置 & 说明\\
\midrule
\srcpath{compute\_carbon\_analysis(sys, pf\_result, opt)} & carbon\_analysis.hpp & 由已收敛潮流结果做快照追踪+矩阵碳核算\\
\srcpath{compute\_carbon\_analysis(sys, pf\_opt, ca\_opt)} & 同上 & 便捷重载：先跑潮流再核算\\
\srcpath{run\_carbon\_analysis(sys, opt)} & 同上 & 向后兼容包装（默认潮流后核算）\\
\srcpath{compute\_annual\_carbon\_analysis(...)} & annual\_carbon\_analysis.hpp & 年度碳（静态/逐步系统/\fld{TimeSeriesPFResult} 四重载）\\
\srcpath{compute\_annual\_user\_gec\_accounting(...)} & 同上 & 用户级绿证净排放核算\\
\srcpath{compute\_annual\_node\_gec\_accounting(...)} & 同上 & 母线级绿证核算\\
\bottomrule
\end{tabular}
\end{center}

\subsection{关键数据结构}
\compmeta{CarbonAnalysisResult}{include/hacdcpf/carbon\_analysis/carbon\_analysis.hpp}
主要字段：\fld{carbon\_sources}、\fld{load\_carbon}/\fld{dc\_load\_carbon}、
\fld{branch\_carbon}/\fld{dc\_branch\_carbon}、\fld{bus\_carbon}/\fld{dc\_bus\_carbon}、
\fld{vsc\_carbon}、\fld{storage\_carbon}，以及有效性字段 \fld{tracing\_verified}、
\fld{matrix\_solved}、\fld{matrix\_rank}、\fld{matrix\_relative\_residual}、
\fld{power\_balance\_verified}。\texttt{EmissionsSummary} 汇总发电/负荷/损耗排放与
平衡误差 \fld{balance\_error\_tco2}/\fld{balance\_error\_pct}。

\section{核心数学模型}
碳流分析建立在已收敛潮流的支路有功流向之上，把每个碳源（发电机、外部电网、
可再生、储能作为碳汇/源、换流器与能量路由器）的排放沿潮流方向分摊到负荷与
支路损耗，AC/DC 集合分开处理。

\subsection{比例功率追踪法（BFS）}
对母线 $i$，令流入功率来自集合 $\mathcal{I}(i)$，其中源母线 $k$ 的贡献占比为
$P_{ki}/P_i^{\text{in}}$，则母线碳强度 $\rho_i$ 满足比例混合律
\begin{equation}
\rho_i = \frac{\sum_{k\in\mathcal{I}(i)} P_{ki}\,\rho_k + \sum_{g\in i} P_{g}\,e_{g}}
              {\sum_{k\in\mathcal{I}(i)} P_{ki} + \sum_{g\in i} P_{g}},
\end{equation}
其中 $e_g$ 为机组排放强度。负荷排放 $=\rho_i P_{d,i}$；支路损耗排放全额
计入送端节点（即 $C_{es}^{loss}=P_e^{loss}F_{f(e),s}$，
\srcpath{proportional\_tracing}，src/carbon\_analysis/carbon\_analysis.cpp:proportional\_tracing），不经
\fld{loss\_allocation\_alpha} 分摊——该系数仅作用于下节的矩阵法
（\srcpath{solve\_carbon\_matrix}，src/carbon\_analysis/carbon\_analysis.cpp:solve\_carbon\_matrix）。

\subsection{碳排放流矩阵法}
在标准碳排放流理论（Kang 等）下，节点碳势向量 $\mathbf{w}$ 满足稀疏线性方程
\begin{equation}
\mathbf{A}\,\mathbf{w} = \mathbf{b},
\end{equation}
其中 $\mathbf{A}$ 为节点入流/支路分配算子、$\mathbf{b}$ 为发电碳注入。实现不构造
稠密 $n\times n$ 矩阵：追踪势用稀疏 LU 一次分解，节点势用秩揭示的稀疏 QR
（\srcpath{SparseQR}）求解，从而对奇异/病态情形保持可见诊断
（\fld{matrix\_rank}、\fld{matrix\_relative\_residual}）。

\section{接口与参数}
\begin{paramtable}
\fld{min\_contribution\_mw} & $P_{\min}$ & MW & $0\sim1$ & — & 追踪贡献下限，低于则忽略\\
\fld{loss\_allocation\_alpha} & $\alpha$ & — & $0\sim1$ & $0.5$ & 损耗排放送/受端分摊系数（仅矩阵法与损耗强度后处理；BFS 追踪恒归送端）\\
\fld{regularization\_eps} & $\varepsilon$ & — & 小正数 & — & 矩阵正则化项\\
\fld{max\_matrix\_condition\_estimate} & — & — & — & — & 条件数上限（超限告警）\\
\fld{verification\_tol} & — & — & 小正数 & — & 平衡校核容差\\
\fld{max\_tracing\_depth} & — & — & 整数 & — & BFS 追踪最大深度\\
\end{paramtable}
年度接口 \srcpath{compute\_annual\_carbon\_analysis} 逐步累计排放，非收敛步的小时
矩阵以 NaN 占位；动态储能碳需为每个潮流步提供一份 \fld{pf\_system\_snapshots[t]}。

\section{验证与边界局限}
\subsection{验证与校核}
注册测试：\srcpath{test\_carbonflow\_tracing}、\srcpath{test\_carbonflow\_dynamic\_storage}、
\srcpath{test\_carbonflow\_case\_validation}，覆盖追踪一致性、动态储能碳与算例校核。

\subsection{边界与局限（如实标注）}
\begin{itemize}
  \item 纯后处理：不重新调度、不修复失败潮流；非收敛输入 $\to$ 默认非验证输出；
  \item 无单一隐式成功标志，须逐项检查有效性字段；
  \item 静态年度重载仅在调度/负荷/储能状态恒定时精确，动态储能碳须逐步快照；
  \item \fld{loss\_allocation\_alpha} 重分配的是责任而非物理损耗。
\end{itemize}
\section{源码等价数学模型}
本章只转写当前实现，不给出比代码更一般的碳流模型。行号对应本章编写时的工作树；函数名是长期定位锚点。

\subsection{年度逐时材料化与 GEC 完整性准入}
\srcpath{src/carbon\_analysis/annual\_carbon\_analysis.cpp:compute\_annual\_carbon\_analysis\_impl}
按首次出现顺序登记各域母线及负荷稳定 ID。每个时段只增长本行，全部处理完后由
\srcpath{extend\_hourly\_width} 各扫描一次三类逐时表，填充后出现列的 NaN。
对 $T$ 时段、$B$ 母线和 $L$ 负荷，材料化工作为 $O(T(B+L))$；旧实现的
$O(T^2(B+2L))$ 全表元数据扫描已移除。失败行和缺失列仍为 NaN；这不改变碳流方程或储能库存递推。

用户和节点 GEC 入口共同调用 \srcpath{validate\_gec\_hourly\_inputs}：时段数和
时长必须为正且有限；逐时能量（MWh）和排放（tCO$_2$）必须严格为
$T\times L$，且每对样本均有限、非负。若提供 \fld{step\_results}，其长度须为 $T$，
每步 \fld{pf\_converged} 与 \fld{carbon\_verified} 均须为真。
不完整输入抛 \texttt{invalid\_argument}，不把未知排放转换成零并报告全年覆盖率。
无诊断的外部组装数据仍可按完整有限矩阵准入，此时不额外声称追踪已被独立验证。
两个 GEC 结果校验器对摘要和逐时指标均先检查有限性；关闭逐时输出不能绕过摘要检查。
注册测试 \srcpath{tests/test\_review\_regressions.cpp} 覆盖上述边界；可复现计时和
构建范围见 \srcpath{docs/testing/module\_code\_audit.md} 的 AUD-091/094/095。

\subsection{节点净源需求与校核}
对应 \srcpath{src/carbon\_analysis/carbon\_analysis.cpp:compute\_node\_net\_source\_requirement}
（1253--1280 行）。对节点 $i$，代码计算
\begin{equation}
R_i=\sum_{l:n(l)=i,\,P_l>kTol}P_l
+\sum_{e:f(e)=i}\max(P_e^{send},0)
-\sum_{e:t(e)=i}\max(P_e^{recv},0).
\end{equation}
若 $|R_i|<10^{-7}$，代码把 $R_i$ 置零。函数参数
\fld{loss\_allocation\_alpha} 在当前函数体中被显式 \texttt{(void)} 忽略；因此此处不得写入损耗分摊项。
校核函数 \srcpath{verify\_node\_power\_balance}（1396--1438 行）计算
\begin{align}
G_i&=\sum_{g:n(g)=i}\max(P_g,0),\\
e_i&=G_i-R_i,\\
\epsilon_P&=\max\left(10^{-7},\ \fld{verification\_tol}\cdot
\max\left(\sum_g\max(P_g,0)+\sum_l\max(P_l,0),1\right)\right),
\end{align}
并仅在 $\max_i|e_i|\le\epsilon_P$ 时令 \fld{power\_balance\_verified}=true。

\subsection{比例追踪稀疏系统}
对应 \srcpath{proportional\_tracing}（1441--1540 行）。代码先定义
\begin{equation}
T_i=\sum_{g:n(g)=i,\,P_g>kTol}P_g
+\sum_{e:t(e)=i}\max(P_e^{recv},0).
\end{equation}
分配矩阵 $M$ 和源注入矩阵 $B$ 的实际装配为
\begin{align}
M_{ii}&=\begin{cases}T_i,&T_i>kTol,\\1,&\text{否则，}\end{cases}\\
M_{t(e),f(e)}&\mathrel{+}=-\max(P_e^{recv},0),\\
B_{i,s}&\mathrel{+}=P_s\quad\text{仅当 }P_s>kTol\text{ 且 }n(s)=i.
\end{align}
使用 \texttt{SparseLU<COLAMD>} 求 $MF=B$。分解或求解失败、$F$ 含非有限值时直接返回未填充结果。
只把 $-10^{-10}<F_{is}<0$ 的微小负数截成零；更大的负数不会在这里自动修正。负荷与损耗来源归属为
\begin{align}
C_{ls}&=P_lF_{n(l),s},\\
C_{es}^{loss}&=P_e^{loss}F_{f(e),s},
\end{align}
且仅当贡献不小于 \fld{min\_contribution\_mw} 才写入结果。

\subsection{碳势矩阵的实际装配}
对应 \srcpath{solve\_carbon\_matrix}（1707--1870 行）。对有效边 $e=(i,j)$，代码先逐项截断
\begin{align}
s_e&=\max(P_e^{send},0),&r_e&=\max(P_e^{recv},0),\\
\ell_e&=\operatorname{clamp}(\max(P_e^{loss},0),0,s_e),&
\alpha&=\operatorname{clamp}(\fld{loss\_allocation\_alpha},0,1),\\
q_e^{recv}&=(1-\alpha)\ell_e,&
q_e^{transfer}&=r_e+\alpha\ell_e,\\
q_e^{send}&=s_e+(2\alpha-1)\ell_e.
\end{align}
随后严格按代码累加
\begin{align}
P_i^{out}&\mathrel{+}=\max(q_e^{send},0),&
P_j^{out}&\mathrel{+}=q_e^{recv},\\
K_{ji}&\mathrel{+}=q_e^{transfer}\quad(q_e^{transfer}>kTol),&
P_{n(l)}^{out}&\mathrel{+}=P_l\quad(P_l>kTol).
\end{align}
右端项为 $b_i=\sum_{g:n(g)=i}\gamma_gP_g$，这里对 $P_g$ 不做 \texttt{max}。普通行的矩阵为
\begin{align}
A_{ii}&=\begin{cases}P_i^{out},&|P_i^{out}|>\epsilon_r,\\\epsilon_r,&\text{否则，}\end{cases}\\
A_{ji}&\mathrel{+}=-K_{ji}.
\end{align}
若一行同时满足负荷小于 $10^{-8}$、发电小于 $10^{-8}$ 且该行最大装配量小于 $10^{-6}$，代码把该行
替换为 $A_{ii}=1,b_i=0$，并不写入该行的转移非对角项。

\subsection{缩放、秩、条件估计与收敛判据}
令 $d_i$ 为装配期间记录的行最大绝对量。代码仅在 $d_i>\max(\epsilon_r,10^{-15})$ 时执行
\begin{equation}
\widetilde A_{ij}=A_{ij}/d_i,\qquad \widetilde b_i=b_i/d_i.
\end{equation}
\texttt{SparseQR<COLAMD>} 的 pivot threshold 为 $\max(\epsilon_r,10^{-12})$。必须满足秩等于节点数，
并以 $R$ 对角枢轴估计
\begin{equation}
\widehat\kappa=\frac{\max_i|R_{ii}|}{\min_i|R_{ii}|}
\end{equation}
有限且不超过 \fld{max\_matrix\_condition\_estimate}（当该上限大于零）。解中若
$w_i<-\max(\epsilon_r,10^{-10})$ 立即失败；其余小负数置零。最终判据精确为
\begin{align}
r&=Aw-b,\\
\rho&=\frac{\|r\|_2}{\max(\|A\|_F\|w\|_2+\|b\|_2,10^{-15})},\\
solved&=\operatorname{finite}(\|r\|_2,\rho)\land
\neg(\|r\|_2>\epsilon_r\land\rho>\fld{verification\_tol}).
\end{align}
注意逻辑连接是“绝对残差和相对残差同时超限才失败”，不是两者任一超限即失败。

\subsection{损耗强度与储能快照结果}
\srcpath{edge\_loss\_carbon\_intensity}（132--140 行）在 $P_e^{loss}>kTol$ 时返回
\begin{equation}
w_e^{loss}=\alpha w_{f(e)}+(1-\alpha)w_{t(e)},
\end{equation}
否则返回零。单步储能结果 \srcpath{build\_storage\_carbon\_result}（1654--1675 行）使用
\begin{align}
E_s&=\max(SOC_s,0)\max(E_s^{rated},0),\\
(w_s,E_s^{CO2})&=\begin{cases}
(w_{charging\_load},E_{charging\_load}^{CO2}),&P_s<-kTol\text{ 且存在充电负荷，}\\
(w_s^{inventory},\max(P_s,0)w_s^{inventory}),&\text{否则。}
\end{cases}
\end{align}
此处 \fld{total\_emissions\_tco2} 数值按 1 小时快照约定生成，函数体没有额外乘时长。

\subsection{年度储能终端快照更新}
对应 \srcpath{carry\_storage\_carbon\_state\_from\_terminal\_snapshot}
（649--728 行）。实现定义
\begin{align}
\eta_c&=\operatorname{clamp}(\eta_c^{raw}>0?\eta_c^{raw}:1,10^{-6},1),\\
\eta_d&=\operatorname{clamp}(\eta_d^{raw}>0?\eta_d^{raw}:1,10^{-6},1),\\
\sigma&=\operatorname{clamp}(1-\max(sd,0)/100,0,1),\\
E_c&=\max(-P_s,0)\Delta t,&E_c^{stored}&=E_c\eta_c,\\
E_d&=\max(P_s,0)\Delta t,&E_d^{withdrawn}&=E_d/\eta_d,\\
E_{sd}&=E_0(1-\sigma),&
\widehat E_1&=E_0\sigma+E_c^{stored}-E_d^{withdrawn}.
\end{align}
能量误差保存为 $\widehat E_1-E_1$。碳损失和终端库存严格为
\begin{align}
C_c^{loss}&=(E_c-E_c^{stored})w_{bus},\\
C_d^{loss}&=(E_d^{withdrawn}-E_d)w_0,\\
C_{sd}^{loss}&=E_{sd}w_0,\\
C_1&=E_0w_0+E_cw_{bus}-E_dw_0-C_c^{loss}-C_d^{loss}-C_{sd}^{loss},\\
w_1&=\begin{cases}\max(C_1,0)/E_1,&E_1>kTol,\\0,&\text{否则。}\end{cases}
\end{align}
年度库存残差对应 \srcpath{finalize\_storage\_adjusted\_balance}（490--534 行）：
\begin{equation}
r_{inv}=C_0+C_c-C_d-(C_c^{loss}+C_d^{loss}+C_{sd}^{loss})-C_T.
\end{equation}

\subsection{复杂度与未实现语义}
构图与装配为 $O(n+m+g+l)$；比例追踪一次稀疏 LU 后求多右端，碳势为一次稀疏 QR，实际复杂度取决于
消元填充。\fld{max\_tracing\_depth} 当前不出现在上述矩阵追踪实现中，不能据字段名宣称存在深度截断。
任何分解失败、秩亏、条件估计超限、显著负解或残差判据失败均保持 \fld{matrix\_solved}=false。

% theory_carbon_flow.tex -- 碳流通用理论层
\section{碳流追踪、矩阵碳势与年度库存理论}\label{sec:carbon-theory}

\begin{theorynote}
本章给出比例追踪、碳势矩阵、损耗分摊和储能库存的通用理论。实现只覆盖本章标注的
有限模型范围；非线性潮流、生命周期制造碳和外部碳因子不因理论公式而自动获得支持。
\end{theorynote}

\subsection{比例追踪}
对有向支路 $e:i\to j$，令 $F_e>0$ 表示按潮流方向的有功流，母线 $j$ 的总入流为
$P_j^{in}$。比例追踪假设来自各上游源的混合比例在母线内保持一致：
\begin{equation}
\rho_j=\frac{\sum_{e\in\delta^-(j)}F_e\rho_{o(e)}+\sum_{g\in j}P_ge_g}{P_j^{in}},
\qquad C_{d,j}=P_{d,j}\rho_j.
\end{equation}
支路损耗若归送端，则 $C_e^{loss}=P_e^{loss}\rho_{i}$；这是一种责任分摊规则，不是
损耗物理量的改变。

\subsection{矩阵碳势}
将节点入流比例写成稀疏矩阵 $A$，源碳注入写成 $b$，节点碳势满足
\begin{equation}Aw=b.
\end{equation}
数值解必须同时检查秩、条件估计和归一化残差
$r=\|Aw-b\|_2/\max(1,\|b\|_2)$。奇异系统不能用零向量代替未知碳势；应返回
\fld{matrix\_solved=false} 或显式的非验证状态。

\subsection{年度储能库存}
储能放电碳强度是库存状态变量而非永久常数：
\begin{align}
E_{t+1}&=\rho E_t+\eta_{ch}P_t^{ch}\Delta t-P_t^{dis}\Delta t/\eta_{dis},\\
C_{t+1}^{inv}&=\rho_C C_t^{inv}+C_t^{ch}-C_t^{dis}.
\end{align}
年度累计排放应按时间顺序更新 $C_t^{inv}$；将某一时刻的
\fld{soc\_carbon\_intensity\_tco2\_mwh} 复制到全年会破坏库存守恒。

\begin{gapnote}
当前模块是已收敛 PF 的后处理；不重新调度、不修复失败潮流，不含制造碳、网络材料碳和
外部生命周期数据库。AC 与 DC 结果向量也不能通过裸 bus ID 合并。
\end{gapnote}

\subsection{与实现的对应}
\begin{center}\footnotesize
\begin{longtable}{@{}L{7.0cm}L{8.2cm}@{}}
\toprule 理论条目 & 实现状态\\\midrule
\endfirsthead\toprule 理论条目 & 实现状态\\\midrule\endhead
比例追踪 & \implfull{src/carbon_analysis/carbon_analysis.cpp:proportional_tracing}\\
稀疏矩阵碳势 & \implfull{src/carbon_analysis/carbon_analysis.cpp:solve_carbon_matrix}\\
节点功率平衡重构 & \implfull{src/carbon_analysis/carbon_analysis.cpp:verify_node_power_balance}\\
逐步储能库存碳 & \implfull{src/carbon_analysis/annual_carbon_analysis.cpp:update_storage_carbon_state}\\
制造/网络材料生命周期碳 & \implnone\\
\bottomrule
\end{longtable}
\end{center}

% numerical_cross_validation.tex -- 碳流数值证据
\section{数值结果与独立方程校核}\label{sec:carbon-numerical}

本轮 macOS Release 直接运行三个碳流目标：
\srcpath{test\_carbonflow\_tracing} 通过 21 cases/246 assertions，
\srcpath{test\_carbonflow\_dynamic\_storage} 通过 15 cases/126 assertions，
\srcpath{test\_carbonflow\_case\_validation} 通过 5 cases/202 assertions。

\subsection{快照追踪证据}
测试覆盖比例追踪、矩阵法、AC/DC 分域、支路损耗责任、零功率源、奇异矩阵诊断和 JSON
结果语义。准入条件为：已收敛 PF 才允许 \fld{tracing\_verified=true}；矩阵法必须报告
有限的秩和相对残差；节点功率平衡误差必须不超过选项 \fld{verification\_tol}。
这些测试是独立于生产调用顺序的闭式/构造网络回归，但没有与 OpenDSS 或第三方碳流软件
进行交叉引擎比较。

\subsection{动态储能证据}
15 个动态储能用例覆盖充电碳汇、放电碳源、端末 SOC 和逐小时快照。测试运行中出现的
“bus has load but no in-service connection”只是构造案例的拓扑 warning，测试仍以结构化
结果和 126 条断言通过；它不能被解释成一般孤岛系统均可验证。

\subsection{独立方程 oracle 交叉验证}\label{sec:carbon-xref}
注册测试 \srcpath{carbon\_analysis\_cross\_validation} 以确定性算例把生产碳流结果与一个
不链接 hacdcpf 的独立 Python oracle 对照。证据生成器
\srcpath{tools/carbon\_analysis\_validation/validate\_carbon\_xref.cpp:emit\_case} 在三个
手工构造、解析可知的交流算例上运行
\srcpath{src/carbon\_analysis/carbon\_analysis.cpp:compute\_carbon\_analysis}，导出比例碳
排放流的输入（机组排放因子与出力、负荷、定向支路潮流）与求解得到的节点
碳强度向量；oracle \srcpath{tools/carbon\_analysis\_validation/run\_cross\_validation.py} 依据
Kang 碳排放流方程独立重建线性系统 $Aw=b$（实现见
\srcpath{src/carbon\_analysis/carbon\_analysis.cpp:solve\_carbon\_matrix}），用独立高斯消元
重解，并三重校核：解析闭式、独立重解 $|w_{oracle}-w_{cpp}|$，以及节点守恒
$\|Aw_{cpp}-b\|$、支路损耗归属 $\rho=\alpha w_{from}+(1-\alpha)w_{to}$ 与系统排放平衡。

复现命令为
\begin{verbatim}
ctest --test-dir build/macos-release \
  -R '^carbon_analysis_cross_validation$' --output-on-failure
\end{verbatim}

三个算例与解析期望为：单源串馈（全节点碳强度等于唯一机组排放因子 $0.5$）、
无损双源混合（负荷母线为出力加权平均 $(30\times1.0+70\times0.0)/100=0.3$）、有损双源
混合（$\alpha=0.5$ 时负荷母线为 $57.5/95$，两条支路损耗强度分别为 $0.5\,w_{from}+0.5\,w_{to}$，
系统排放平衡闭合到 $60$ tCO$_2$）。准入阈值为 $10^{-7}$；本轮全部七类检查的最大
误差为 $3.553\times10^{-15}$，负控制（人为扰动一个导出节点强度）如期触发失败。

\subsection{数值边界}
矩阵秩亏、潮流不收敛、非有限输入和缺少逐时系统快照会使结果进入未验证路径；报告必须
同时保留 \fld{matrix\_solved}、\fld{power\_balance\_verified} 和错误列表。第~\ref{sec:carbon-xref}~节
的独立方程 oracle 现以机器可读结果证明交流比例碳流矩阵解与已发布方程一致，但它重解同一线性
系统，属独立实现校核而非第二套物理模型；VSC/DC-DC/储能/外送碳路径以及与第三方碳流
软件的跨引擎比较仍未覆盖，故除该 oracle 声明的交流范围外，本节不声明碳流估计的
外部准确度。

\section{跨模块工业级技术评价基线}

本章规定本手册所述模块进入工程算例、批量研究或生产服务前必须共同满足的评价基线。
具体模型公式、变量、约束和专用算法以正文相应章节为准；本章不扩大源码已经实现的范围。

\subsection{输入、输出、边界条件与数据结构审计}
输入审计应逐项检查有限性、单位、上下界、枚举值和引用完整性。系统对象必须先按所需层级完成
校验；AC 与 DC 母线采用域限定映射，组件稳定 \texttt{index}、向量位置和临时图索引不得混用。
输出审计必须记录单位、索引空间、模型范围、收敛或可行状态、后端、容差、迭代/节点计数和告警。
空系统、空时域、孤岛、重复 ID、零容量、退化约束与不可用可选后端均属于边界测试集。

\subsection{工程假设与数值稳定性分析}
模型使用的平衡、线性化、稳态、独立性、概率分布、时间离散或网络等值假设必须与应用场景匹配。
数值评价至少包含变量缩放、绝对/相对残差、可行性违反、条件数或枢轴质量、步长/阻尼、迭代停滞
和随机种子复现性。若采用正则化、平滑、回退、松弛或近似，应同时报告触发条件、参数值、对主指标
的敏感性以及是否改变模型口径；不得用“已返回结果”替代收敛或有效性证明。

\subsection{复杂度与资源分析}
令 $n_x$、$n_c$、$n_t$、$n_s$、$|V|$、$|E|$ 分别表示变量、约束、时段、场景、节点和边数。
报告应分别给出建模、求解、后处理的时间与内存。图遍历通常为 $O(|V|+|E|)$；逐时段/逐场景
流程至少线性增长为 $O(n_t n_s)$ 次子问题调用；稠密线性代数最坏可达 $O(n_x^3)$，稀疏方法则应
报告结构非零数、填充率和因子分解峰值内存；MILP/NLP 不给出虚假的多项式最坏界，而报告节点数、
间隙、时限和实测扩展曲线。复杂度结论必须基于固定硬件、固定后端和可复现算例族。

\subsection{异常工况处理}
非法输入应在进入求解器前返回结构化错误；不可行、无界、不收敛、时限、内存不足和后端缺失必须
相互区分。部分结果只有在对应 \fld{ValidityFlags}、\fld{model\_scope}、
\fld{model\_limitations} 或等价字段允许时才可使用。异常路径不得静默丢弃 AC/DC 资产、制造平衡
节点、把 NaN 改写为零或把近似解标为全局最优；系统原始对象应保持可恢复和可重试。

\subsection{与其他模块的接口关系}
接口链路遵循“工程语义模型 $\to$ 校验 $\to$ 投影/装配 $\to$ 求解或分析 $\to$ 结果回投影”。
跨模块传递时保留稳定 ID、单位、符号、时间基准和场景身份；消费潮流、OPF、动态或可靠性结果前，
先检查其收敛、覆盖范围和有效性标志。HTTP/JSON/Python 层只做语义保持适配，不重新解释求解结果。

\subsection{测试与验证建议}
最低验证矩阵包括：手算小例、标准算例、边界/异常输入、随机性质测试、算法或后端交叉校核、
序列化往返、稳定 ID 回投影、AC/DC 同号母线、重复运行确定性和规模扩展测试。数值算法还应加入
残差独立重算、雅可比有限差分或自动微分校核；近似模型应与高保真模型比较并给出误差分布，而非
只报告单个平均值。回归报告必须列明构建配置、算例版本、容差、通过/失败/跳过数及未验证范围。

\subsection{工业级评估指标}
\begin{itemize}
  \item \textbf{正确性}：守恒残差、约束最大违反、解析/基准误差、结果回投影一致性；
  \item \textbf{鲁棒性}：收敛率、不可行识别率、异常分类准确率、扰动敏感性与随机复现率；
  \item \textbf{性能}：端到端时延、吞吐量、峰值内存、并行效率与规模增长斜率；
  \item \textbf{可追溯性}：输入模式版本、稳定 ID 覆盖率、模型范围/限制字段完整率、告警可定位率；
  \item \textbf{工程有效性}：与标准、外部引擎或现场数据的偏差，以及误判/漏判对业务指标的影响。
\end{itemize}
指标应预先给出验收阈值和失败处置；未达到阈值时只能标记为实验性或受限适用。

\subsection{适用范围、局限性与后续改进}
本手册只评价当前源码覆盖的模型、后端和数据结构，不对未建模物理过程、未安装求解器或未执行的
外部验证作保证。后续改进应按“缺口 $\to$ 可量化指标 $\to$ 源码/测试变更 $\to$ 独立验证”闭环，
优先消除会造成错误结论的语义缺口，其次提升数值鲁棒性与性能，最后扩展模型范围。


\end{document}
